\clearpage
\VerdierPhysicalPageAnchor{0068}%
\VerdierSemanticAnchor{ch1-2-2-1-continuation}%
$j<i$ 时 $\varepsilon(e_j,i)=-1$，$i\leq j$ 时
$\varepsilon(e_j,i)=1$。而且，映射
$(\sigma,i)\mapsto\varepsilon(\sigma,i)$ 由这些性质唯一确定。由此立即得到关系

\VerdierSemanticAnchor{ch1-2-2-1-eq2}
\begin{equation}
  \varepsilon(\sigma,i)\varepsilon(\sigma+e_i,j)
  =-\varepsilon(\sigma,j)\varepsilon(\sigma+e_j,i)
  \tag{2.2.1.2}
\end{equation}
它对每个 $\sigma\in\mathbb Z[n]$ 与每一对
$(i,j)\in[n]\times[n]$，$i\neq j$ 都成立。

\VerdierSemanticAnchor{ch1-2-2-2}
\noindent\textbf{2.2.2.}\quad
设 $Y$ 为 $C$ 中的一个 $n$-重复形。令

\VerdierSemanticAnchor{ch1-2-2-2-eq1}
\begin{equation}
  \delta_Y^{\sigma,i}=\varepsilon(\sigma,i)d_Y^{\sigma,i}.
  \tag{2.2.2.1}
\end{equation}

\VerdierSemanticAnchor{ch1-2-2-2-eq2}
由 (2.2.1.2)，关系 (2.1.2.1) 等价于关系
\begin{equation}
  \begin{aligned}
    \delta_Y^{\sigma+e_i,i}\delta_Y^{\sigma,i}
      &=0,
      &&\forall\sigma\in\mathbb Z[n],\quad\forall i\in[n],\\
    \delta_Y^{\sigma+e_j,i}\delta_Y^{\sigma,j}
      +\delta_Y^{\sigma+e_i,j}\delta_Y^{\sigma,i}
      &=0,
      &&\forall\sigma\in\mathbb Z[n],\quad\forall i,j\in[n].
  \end{aligned}
  \tag{2.2.2.2}
\end{equation}

\VerdierSemanticAnchor{ch1-2-2-2-pushforward}
现设 $\varphi\colon[n]\to[m]$ 为映射，并设 $Y$ 是
$\varphi$-可求和的。我们用下列公式定义一个 $m$-重复形 $\int_\varphi Y$：

\VerdierSemanticAnchor{ch1-2-2-2-eq3}
\begin{equation}
  \left(\int_\varphi Y\right)^\tau
  =\bigoplus_{\varphi(\sigma)=\tau}Y^\sigma,
  \qquad \forall\tau\in\mathbb Z[m],
  \tag{2.2.2.3}
\end{equation}

\VerdierSemanticAnchor{ch1-2-2-2-eq4}
\begin{equation}
  \begin{aligned}
    \delta_{\int_\varphi Y}^{\tau,i}(y_\sigma)
      &=\sum_{\varphi(j)=i}\delta_Y^{\sigma,j}(y_\sigma),\\[-1mm]
    &\hspace{-18mm}
      \sigma\in\mathbb Z[n],\quad i\in[m],\quad
      y_\sigma\in Y^\sigma,\quad\tau=\varphi(\sigma).
  \end{aligned}
  \tag{2.2.2.4}
\end{equation}

\VerdierSemanticAnchor{ch1-2-2-2-additive-interpretation}
当 $C$ 是阿贝尔群范畴时，公式 (2.2.2.4) 定义了态射
\[
  \left(\int_\varphi Y\right)^\tau
  =\bigoplus_{\varphi(\sigma)=\tau}Y^\sigma
  \longrightarrow
  \left(\int_\varphi Y\right)^{\tau+e_i}
  =\bigoplus_{\varphi(\sigma)=\tau+e_i}Y^\sigma.
\]
当 $C$ 是加法范畴时，
它也能定义从 $\left(\int_\varphi Y\right)^\tau$ 到
$\left(\int_\varphi Y\right)^{\tau+e_i}$ 的态射：只需在 $C$ 上取值于阿贝尔群
范畴的反变函子范畴中工作，即把 $y_\sigma$ 解释为从 $C$ 中任意对象到
$Y^\sigma$ 的态射。

\VerdierSemanticAnchor{ch1-2-2-2-verification}
直接计算可知 $\delta_{\int_\varphi Y}^{\tau,i}$ 满足关系 (2.2.2.2)。因此回到
微分后，由 (2.2.2.3) 定义的对象配以下列态射：

\VerdierSemanticAnchor{ch1-2-2-2-eq5}
\begin{equation}
  \begin{aligned}
    d_{\int_\varphi Y}^{\tau,i}(y_\sigma)
      &=\varepsilon(\tau,i)
        \sum_{\varphi(j)=i}
          \varepsilon(\sigma,j)d_Y^{\sigma,j}(y_\sigma),\\[-1mm]
    &\hspace{-18mm}
      \tau\in\mathbb Z[m],\quad i\in[m],\quad
      \sigma\in\mathbb Z[n],\quad\varphi(\sigma)=\tau
  \end{aligned}
  \tag{2.2.2.5}
\end{equation}
